\documentclass[a4paper,12pt]{article}
\usepackage[utf8]{inputenc}
\usepackage[T1]{fontenc}
\usepackage[margin=2.5cm]{geometry}
\usepackage{ctex} % 中文支持
\usepackage{xcolor}
\usepackage{titlesec}
\usepackage{setspace}

% -----------------------------------------------------------
% 格式设置区：Level 3 纯文本墙 + Level 2 专业语态 + Level 2 合规披露
% -----------------------------------------------------------

% 仅保留一级标题，格式简洁，强调文本的连续性
\titleformat{\section}{\Large\bfseries\color{blue!40!black}}{\thesection}{1em}{}

% 设置行间距，在密集文本中保持基本的阅读区分度
\onehalfspacing
\setlength{\parskip}{1em}
\setlength{\parindent}{2em}

% -----------------------------------------------------------
% 文档元数据
% -----------------------------------------------------------
\title{\textbf{2024年度环境、社会及管治（ESG）报告：可持续发展管理与绩效}}
\author{本公司董事会办公室}
\date{2025年3月}

\begin{document}

\maketitle

\section{前言与环境合规管理体系}

2024年，本公司坚持将环境、社会及管治（ESG）理念深度融入企业发展战略与日常运营管理体系之中。作为国内领先的化工行业上市公司，我们严格遵守《环境保护法》、《安全生产法》及相关监管要求，致力于通过技术创新与精细化管理，实现经济效益与社会效益的协同发展。在环境管理方面，公司建立了完善的环境管理体系，坚持源头减量、过程控制与末端治理相结合的原则，切实履行环境保护主体责任。报告期内，公司加强了环境数据的监测与统计，确保披露数据的真实性与准确性。根据年度合规监测数据，2024年公司温室气体排放总量为39,107.65吨CO2e，排放强度控制在4.58吨/百万元营收，体现了公司在低碳转型方面的管理成效。在废弃物管理方面，公司严格执行分类处置制度，全年产生一般工业固废170.48吨；对于环境风险较高的危险废物，合规处置率达到100\%，确保无环境污染风险。在资源利用方面，全年新鲜水总用量为149.55万吨，包装材料使用总量为7,677吨。为保障上述环境合规目标的实现，公司全年投入环保资金457.93万元，并依法缴纳环保税9.6万元，体现了企业对环境外部成本的主动承担。

为积极响应国家“双碳”战略及行业清洁生产要求，公司在生产及运营环节实施了多项技术改造与管理优化措施。在固定源排放管理上，核心生产基地完成了供热系统的清洁化改造，通过引入清洁能源热源替代原有的燃煤锅炉设施，该项目的实施有效消除了燃烧环节的废气排放，实现了生产过程废气的“零排放”，符合国家关于工业锅炉整治的合规要求。此外，旗下化工厂实施了冷凝水回收利用项目，通过回收生产蒸汽冷凝水作为锅炉补水，实现了热能与水资源的双重回收，显著提升了资源利用效率。在移动源排放管理上，针对西北某大型露天煤矿项目的运输环节，公司积极推进运输设备的电气化升级，以降低化石能源消耗。本年度，公司新增21台纯电动矿卡及36台油电增程式矿卡，其中纯电动矿卡具备制动能量回收功能，显著降低了运营能耗；同时，公司前瞻性地投入了12台无人驾驶电动矿卡，通过自动驾驶技术优化行驶路径，实现了作业过程的零排放与低噪音，最大限度减少了对矿区周边生态环境的影响。

\section{研发创新与行业高质量发展}

技术创新是公司核心竞争力的体现，也是推动化工与民爆行业安全、高效发展的关键动力。2024年，公司持续保持高强度的研发投入，并构建了完善的知识产权保护体系，以驱动行业的高质量发展。报告期内，公司研发支出共计4.19亿元，占营业收入比例为4.91\%。公司拥有一支由1,064名专业人员组成的研发队伍，本年度新增授权专利126项，累计拥有有效专利596项。这些投入与产出为公司的可持续发展提供了坚实的技术支撑。公司重点推广了“现场混装水胶炸药工艺”，该技术实现了原材料运输与混合工序的分离，填补了国内技术空白，其应用价值在于通过现场即时混合，消除了成品炸药运输过程中的安全风险，同时具备密度可调特性，能够根据地质条件优化爆破效果。在特种爆破领域，公司自主研发的“二氧化碳相变膨胀激发管”产能规模达500万只，该技术利用液态二氧化碳物理相变原理破岩，全过程无明火产生，显著提升了高瓦斯矿井等复杂环境下的开采安全性。此外，公司构建的“智能矿山作业系统”集成了5G通讯、边缘计算与机器人技术，实现了5G钻机、自动验孔机器人与无人矿卡的协同作业，有效降低了高危岗位的人员暴露风险，推动矿山作业向无人化、智能化转型，符合行业技术进步的导向。

\section{社会责任履行与员工权益保障}

公司坚持“以人为本”的管理理念，切实保障员工权益，并积极投身公益事业，回馈社会。截至2024年末，公司员工总数为7,399人，劳动合同签订率保持100\%。公司始终将安全生产置于首位，构建了全方位的安全保障体系。为转移及化解安全风险，公司投入355.99万元购买安全生产责任险，并投入892.5万元用于工伤保险。同时，公司建立了“1+14+N”的应急预案体系，确保突发状况下的快速响应。在海外运营实践中，纳米比亚矿服项目部克服跨文化管理挑战，建立了标准化的安全管理制度，如严格的晨会制度及“手指口述”确认法。报告期内，该项目实现连续130万小时无损工事件，展现了公司在国际化运营中对员工生命安全的高度负责，符合国际劳工保护标准。

在履行企业公民责任方面，公司积极响应国家乡村振兴战略，并在突发灾害救援中发挥专业优势。2024年，公司对外捐赠及公益支出总额1,219万元，其中乡村振兴专项投入1,125万元。在2024年湖南某地决堤抢险行动中，公司迅速响应，调集专业技术力量。通过采用“松动爆破”技术方案，在确保石料完整性的前提下，短时间内开采石料8.4万吨，为抗洪抢险提供了关键的物资保障。这一行动不仅体现了企业的社会担当，也证明了公司将专业能力转化为社会价值的能力。

\section{公司治理与商业道德建设}

作为上级集团成员企业，公司致力于完善治理结构，维护投资者权益，并打造廉洁透明的供应链。公司构建了“3+3+3”的董事会结构，即3名股东董事、3名独立董事和3名执行董事，这种平衡的架构确保了决策的科学性与制衡性。凭借规范的运作与透明的披露，公司在证券交易所年度信息披露考评中获得“A级”评价。公司重视股东投资回报，2024年度派发现金红利2.54亿元（每10股2.05元），分红比例占可分配利润的40.12\%，切实让投资者分享企业发展红利。在供应链管理方面，公司对1,548家合作供应商实施严格的商业道德管理，推行《廉洁协议》签署制度，签署率达到100\%。这一举措有效防范了商业贿赂风险，确保了采购过程的公开、透明，共同营造健康的行业生态。

\end{document}